\documentclass[11pt]{amsart}

\usepackage[T1]{fontenc}
\usepackage[utf8]{inputenc}
\usepackage{lmodern}
\usepackage{amsmath,amssymb,amsthm}
\usepackage[margin=1.1in]{geometry}
\usepackage{enumitem}
\usepackage[hidelinks]{hyperref}

\theoremstyle{plain}
\newtheorem{theorem}{Theorem}[section]
\newtheorem{proposition}[theorem]{Proposition}
\newtheorem{lemma}[theorem]{Lemma}
\newtheorem{corollary}[theorem]{Corollary}
\newtheorem*{theoremA}{Theorem A}
\theoremstyle{remark}
\newtheorem{remark}[theorem]{Remark}

\newcommand{\Q}{\mathbb{Q}}
\newcommand{\Z}{\mathbb{Z}}
\newcommand{\Qp}{\mathbb{Q}_p}
\newcommand{\Qpnr}{\widehat{\mathbb{Q}}_p^{\mathrm{nr}}}
\newcommand{\Zpnr}{\widehat{\mathbb{Z}}_p^{\mathrm{nr}}}
\newcommand{\Ohat}{\widehat{\Omega}}
\newcommand{\Fr}{\mathrm{Fr}}
\newcommand{\nrd}{\mathrm{nrd}}
\newcommand{\PGL}{\mathrm{PGL}}
\newcommand{\GL}{\mathrm{GL}}
\newcommand{\Aut}{\mathrm{Aut}}

\setlist[itemize]{leftmargin=*}

\title{The special fibre of Atkin--Lehner quotients at non-squarefree level}

\begin{document}

\maketitle

\section{Construction of the dual graph}

Let $B$ be an indefinite quaternion algebra over $\Q$ of discriminant $D > 1$, let $N \ge 1$ be prime to $D$, let $X = X_0^D(N)$, let $W$ be a group of Atkin--Lehner involutions of $X$, and put $Y = X/W$. Fix a prime $p \mid D$. Let $B'$ be the definite quaternion algebra of discriminant $D/p$, and let $R \supset S$ be Eichler orders of $B'$ of levels $N$ and $Np$. For a lattice $L$ in $B'$ and a prime $\ell$ write $L_\ell = L \otimes \Z_\ell$, and for a locally principal left ideal $I$ write $[I]$ for its left ideal class. All ideals below are locally principal.

Since $p \nmid D/p$, the algebra $B'$ splits at $p$. Put $\varpi = \left(\begin{smallmatrix} 0 & 1 \\ p & 0 \end{smallmatrix}\right)$. We fix $B' \otimes \Qp \cong M_2(\Qp)$ with $R_p = M_2(\Z_p)$ and $S_p = \left(\begin{smallmatrix} \Z_p & \Z_p \\ p\Z_p & \Z_p \end{smallmatrix}\right)$. This is possible: $S_p$ is conjugate to this standard order \cite[Proposition~23.4.3(ii)]{Voi}, which is $M_2(\Z_p) \cap \varpi^{-1}M_2(\Z_p)\varpi$ \cite[(23.1.4)]{Voi}; these are the only two maximal orders containing it \cite[23.5.16]{Voi}, so $R_p$ is one of them, and conjugating by $\varpi$, which normalises the standard order, we may take it to be $M_2(\Z_p)$.

We next define the Atkin--Lehner ideals. Let $O$ be $R$ or $S$, and let $m$ be a divisor of the reduced discriminant of $O$ (that is, of $DN/p$ or of $DN$) with $\gcd(m, \text{disc}/m) = 1$. The \emph{Atkin--Lehner ideal} $J_m(O)$ is the two-sided ideal of $O$ with $J_m(O)_\ell = O_\ell$ for $\ell \nmid m$ and, for $\ell \mid m$, writing $\ell^e \parallel m$, $J_m(O)_\ell = O_\ell\, g_\ell$, where $g_\ell$ is an element of the normaliser of $O_\ell$ whose reduced norm has $\ell$-adic valuation $e$, not lying in $\Q_\ell^\times O_\ell^\times$. Such $g_\ell$ exist and all generate the same ideal: for $\ell \nmid D/p$ the normaliser modulo $\Q_\ell^\times O_\ell^\times$ is generated by the class of $g_\ell = \left(\begin{smallmatrix} 0 & 1 \\ \ell^e & 0 \end{smallmatrix}\right)$, in coordinates where $O_\ell$ is upper triangular modulo $\ell^e$ \cite[Proposition~23.4.14]{Voi}, and for $\ell \mid D/p$, $O_\ell$ is the maximal order of a division algebra and $J_m(O)_\ell$ is its maximal ideal \cite[13.3.10, Theorem~13.3.11]{Voi}. In particular $J_p(S)_p = S_p\varpi$. The condition $g_\ell \notin \Q_\ell^\times O_\ell^\times$ matters: when $e \ge 2$ is even, $\ell^{e/2}O$ is another two-sided ideal of the same norm.

Let $\mathcal{G}$ be the graph whose vertices are two copies of the set of classes of locally principal left ideals of $R$, written $([J], 1)$ and $([J], 2)$, and whose edges are the classes of locally principal left ideals of $S$. The edge $[I]$ joins $([RI], 1)$ to $([R\,J_p(S)\,I], 2)$. For $D = pq$ this is the graph of Ribet \cite[Proposition~4.4]{Rib}, whose edges join the images of an ideal class of $S$ under the two Eichler orders of level $N$ containing $S$. For $m \parallel DN$ define maps on the vertices and edges of $\mathcal{G}$ as follows: if $p \nmid m$,
\[
w_m\colon ([J], i) \mapsto ([J_m(R)\,J], i), \qquad [I] \mapsto [J_m(S)\,I];
\]
if $m = p$,
\[
w_p\colon ([J], i) \mapsto ([J], 3 - i), \qquad [I] \mapsto [J_p(S)\,I];
\]
and $w_m = w_p w_{m/p}$ if $p \mid m$. These are well defined on classes because the $J_m$ are two-sided.

\begin{lemma}\label{lem:action}
The maps $w_m$, for $w_m \in W$, define an action of $W$ on $\mathcal{G}$ by graph automorphisms.
\end{lemma}

\begin{proof}
Atkin--Lehner ideals at different primes commute, since locally one of them is the order itself, and $J_m(O)^2 = mO$ \cite[Proposition~23.4.14, Theorem~13.3.11]{Voi}, with $m$ central; so the maps commute and square to the identity on classes. It remains to check that they preserve incidence. For $p \nmid m$ we have $R\,J_m(S) = J_m(R)$, as both equal $R_\ell g_\ell$ at $\ell \mid m$ and $R_\ell$ elsewhere; hence the edge $[J_m(S)\,I]$ joins $([J_m(R)\,RI], 1)$ to $([J_m(R)\,R\,J_p(S)\,I], 2)$, the images of the endpoints of $[I]$. For $w_p$, the edge $[J_p(S)\,I]$ joins $([R\,J_p(S)\,I], 1)$ to $([R\,J_p(S)^2 I], 2) = ([RI], 2)$, since $J_p(S)^2 = pS$; these are the images under $w_p$ of the endpoints $([RI],1)$ and $([R\,J_p(S)\,I],2)$ of $[I]$.
\end{proof} For a group $H$ acting on a graph $\mathcal{T}$, write $(H\backslash\mathcal{T})^*$ for the quotient graph with the images of the edges inverted by $H$ removed.

\begin{theorem}\label{thm:main}
Over $\Z_{p^2}$, the ring of integers of the unramified quadratic extension of $\Qp$, the curve $Y$ has a projective semistable model, not necessarily regular, whose special fibre has rational components with normalisations isomorphic to $\mathbb{P}^1$, and whose dual graph is $(W\backslash\mathcal{G})^*$.
\end{theorem}

For squarefree $N$, \cite[Theorem~2.6]{PS} gives the corresponding statement for the minimal resolution of this model, whose dual graph is obtained from $(W\backslash\mathcal{G})^*$ by replacing the edge of each node of type $A_{m-1}$ with a chain of $m$ edges; it rests on \cite[p.~202]{Ogg} and \cite[Theorem~1.3]{BD}, the latter stated for squarefree level.

\section{Proof of Theorem~\ref{thm:main}}\label{sec:proof}

\subsection{Uniformization}

We use the identification $B' \otimes \Qp \cong M_2(\Qp)$ fixed in Section~1. Let $\Gamma$ be the image of $R[1/p]^\times$ in $\PGL_2(\Qp)$ and $\Gamma_+ \subset \Gamma$ the subgroup of elements whose reduced norm has even $p$-adic valuation; both are discrete and cocompact. Since $B'$ splits at $p$, the reduced norm maps $R[1/p]^\times$ onto the positive units $p^{\Z}$ of $\Z[1/p]$ \cite[Corollary~28.6.4]{Voi}, so $\Gamma_+$ has index two in $\Gamma$. Let $\Omega$ be Drinfeld's $p$-adic upper half plane, $\Ohat$ its formal model over $\Z_p$, and $I$ the Bruhat--Tits tree of $\PGL_2(\Qp)$.

\begin{proposition}[\v{C}erednik, Drinfeld]\label{prop:BC}
There is an isomorphism of rigid spaces $X \otimes \Qpnr \cong \Gamma_+\backslash\Omega$. Moreover $X$ has a projective model over $\Z_{p^2}$ whose formal completion is $(\Gamma_+\backslash\Ohat) \otimes \Z_{p^2}$, and the descent datum to $\Z_p$ is Frobenius composed with any element of $\Gamma - \Gamma_+$.
\end{proposition}

\begin{proof}
This is \cite[III, Theorem~(5.2) and 5.3.2--5.3.3]{BC}, which applies to the Shimura curve of any order of $B$ that is maximal at $p$, in particular to an Eichler order $O_N$ of level $N$. The two hypotheses of \cite[III, 5.3.3]{BC} hold for $O_N$: for every $\ell \neq p$ the reduced norm maps $(O_N \otimes \Z_\ell)^\times$ onto $\Z_\ell^\times$, since for $\ell \nmid D$ the order $O_N \otimes \Z_\ell$ contains $\mathrm{diag}(u, 1)$ for $u \in \Z_\ell^\times$, and for $\ell \mid D/p$ it is the maximal order of a division algebra, whose units have every element of $\Z_\ell^\times$ as reduced norm \cite[Lemma~13.4.9]{Voi}; and the rational numbers that are units in $O_N \otimes \Z_\ell$ for all $\ell \neq p$ are $\pm p^{\Z}$, which map onto $\Z$ under the $p$-adic valuation. In particular the discrete group of \cite[III, 5.3.3]{BC} is the image of $R[1/p]^\times$ in $\PGL_2(\Qp)$, that is, $\Gamma$.
\end{proof}

\subsection{The Atkin--Lehner involutions}

\begin{lemma}\label{lem:wm}
Let $m \parallel DN$ with $p \nmid m$. Then $J_m(R) \cdot R[1/p] = \alpha_m R[1/p]$ for some $\alpha_m \in R[1/p]$ of reduced norm $m$, and $w_m$ acts on $\Gamma_+\backslash\Omega$ through $\alpha_m$. The element $\alpha_m$ normalises $\Gamma$ and $\Gamma_+$.
\end{lemma}

\begin{proof}
This is stated in \cite[p.~202]{Ogg}, at any level. We give a proof. The isomorphism of \cite[III, Theorem~(5.2)]{BC} is compatible with the action of elements of $(B \otimes \Q_\ell)^\times$, for $\ell \neq p$, that normalise $O_N \otimes \Z_\ell$; in particular with the local Atkin--Lehner elements at the primes $\ell \mid m$. When $O_N^\times$ has torsion, \cite[III, 5.3.2]{BC} deduces the uniformization from that of the subgroup of units congruent to $1$ modulo $M$, for some $M \ge 3$ prime to $DNp$; this subgroup is normalised by the Atkin--Lehner elements, so the compatibility passes to the quotient. Under \cite[III, 5.3.3]{BC}, the action of the local Atkin--Lehner elements becomes that of an element of $B'^\times$ generating $J_m(R) \otimes \Z_\ell$ at every $\ell \neq p$, that is, a generator of $J_m(R) \cdot R[1/p]$. Such a generator exists because $R[1/p]$ is an Eichler $\Z[1/p]$-order in an algebra split at $p$, hence has class number one \cite[Theorem~28.5.5]{Voi}, the narrow class group of $\Z[1/p]$ being trivial. Its reduced norm is $m p^k$ for some $k$, and multiplying by a unit of $R[1/p]$ of reduced norm $p^{-k}$, which exists by \cite[Corollary~28.6.4]{Voi}, we may take its reduced norm to be $m$. It normalises $R[1/p]$ because $J_m(R)$ is two-sided.
\end{proof}

\begin{lemma}\label{lem:wp}
The involution $w_p$ acts on $\Gamma_+\backslash\Omega$ through any element of $\Gamma - \Gamma_+$. Consequently, for $m \parallel DN$ with $p \mid m$, $w_m$ acts through $\alpha_{m/p}\,\omega$ for any $\omega \in \Gamma - \Gamma_+$.
\end{lemma}

\begin{proof}
This is \cite[\S4, p.~463]{Rib} for $D = pq$, combining Drinfeld's description of the action of $(B \otimes \Qp)^\times$ with \cite[Theorems~4.3$'$ and~4.4]{JL}; see also \cite[p.~202]{Ogg}. We give a proof at our generality. Let $\widetilde{\Gamma} = R[1/p]^\times \subset \GL_2(\Qp)$, acting on $\Ohat \widehat{\otimes} \Zpnr$ by $g \otimes \Fr^{-v}$, where $v$ is the $p$-adic valuation of $\det g$ \cite[II, Theorem~(9.3)]{BC}; by \cite[III, 5.3.3]{BC}, the formal completion over $\Z_p$ of the model of Proposition~\ref{prop:BC} is $\widetilde{\Gamma}\backslash(\Ohat \widehat{\otimes} \Zpnr)$. The involution $w_p$ is induced by a uniformizer $\pi$ of $B \otimes \Qp$; the uniformization is compatible with the action of $(B \otimes \Qp)^\times$ \cite[III, Theorem~(5.5)]{BC}, applied with $n = 0$, where $(B \otimes \Qp)^\times$ acts as in \cite[II, 9.4]{BC}; and $\pi$ acts on $\Ohat \widehat{\otimes} \Zpnr$ as $1 \otimes \Fr^{-1}$ \cite[II, Theorem~(9.5)]{BC}. Let $\omega \in \widetilde{\Gamma}$ have reduced norm of odd valuation $2k + 1$. The automorphisms $1 \otimes \Fr^{-1}$ and $\omega^{-1} \otimes 1$ of $\Ohat \widehat{\otimes} \Zpnr$ both normalise the action of $\widetilde{\Gamma}$, since conjugation preserves the valuation of the determinant, so both induce automorphisms of the quotient. They differ by $\omega \otimes \Fr^{-1}$, which is the action of $\omega p^{-k} \in \widetilde{\Gamma}$, the scalar $p$ acting as $1 \otimes \Fr^{-2}$; so they induce the same automorphism of the quotient, and this automorphism is $w_p$. By \cite[III, 5.3.3]{BC}, the base change of the quotient to $\Zpnr$ is $(\Gamma_+\backslash\Ohat) \widehat{\otimes} \Zpnr$, on which $\omega^{-1} \otimes 1$ is the automorphism $\omega^{-1}$; and $\omega^{-1} \equiv \omega$ modulo $\Gamma_+$. Neither the level nor the number of primes dividing $D$ enters. The last assertion follows from Lemma~\ref{lem:wm}, since $w_m = w_p w_{m/p}$.
\end{proof}

\subsection{The quotient by $W$}

Let $\Gamma_W \subset \PGL_2(\Qp)$ be the subgroup generated by $\Gamma_+$ and the elements of Lemmas~\ref{lem:wm} and~\ref{lem:wp} attached to the elements of $W$.

\begin{lemma}\label{lem:quot}
The group $\Gamma_W$ is discrete and cocompact, $\Gamma_W/\Gamma_+ \cong W$, and the quotient by $W$ of the model of Proposition~\ref{prop:BC} is a projective model of $Y$ over $\Z_{p^2}$ with formal completion $\Gamma_W\backslash\Ohat$.
\end{lemma}

\begin{proof}
By Lemmas~\ref{lem:wm} and~\ref{lem:wp}, $\Gamma_W$ normalises $\Gamma_+$ and maps onto $W \subset \Aut(\Gamma_+\backslash\Omega)$. An element $\gamma \in \Gamma_W$ acting trivially satisfies $\gamma z \in \Gamma_+ z$ for every $z \in \Omega$; since $\Gamma_+$ is countable and an element of $\PGL_2(\Qp)$ other than the identity has at most two fixed points, $\gamma$ agrees with a single element of $\Gamma_+$. So the kernel is $\Gamma_+$, $\Gamma_W/\Gamma_+ \cong W$, and $\Gamma_W$, containing $\Gamma_+$ with finite index, is discrete and cocompact. The finite group $\Gamma_W/\Gamma_+$ acts on the formal completion $(\Gamma_+\backslash\Ohat) \otimes \Z_{p^2}$ of the projective model of Proposition~\ref{prop:BC}; by Grothendieck's existence theorem \cite[III, 5.4.1]{EGA} these automorphisms come from unique automorphisms of the model, which on the generic fibre are those of $W$ by Lemmas~\ref{lem:wm} and~\ref{lem:wp}. The quotient of a projective scheme by a finite group is projective, and its generic fibre is the quotient of the generic fibre, here $Y$. Formation of $W$-invariants commutes with the flat base change to the completion, so the completion of the quotient is $(\Gamma_+\backslash\Ohat)/(\Gamma_W/\Gamma_+) = \Gamma_W\backslash\Ohat$, as in Kurihara's construction \cite[p.~285]{Kur}.
\end{proof}

\begin{proposition}[Kurihara]\label{prop:kur}
Let $\Gamma' \subset \PGL_2(\Qp)$ be a discrete cocompact subgroup, possibly with torsion and possibly containing elements that invert edges of $I$. Then $\Gamma'\backslash\Ohat$ is a semistable formal scheme over $\Z_p$. The components of its special fibre are rational curves with normalisations isomorphic to $\mathbb{P}^1$, in bijection with the vertices of $(\Gamma'\backslash I)^*$, and its singular points are in bijection with the edges of $(\Gamma'\backslash I)^*$; the singularity at an edge whose stabiliser has order $m$ is of type $A_{m-1}$.
\end{proposition}

\begin{proof}
This is \cite[\S3, Propositions~3-2 and~3-4]{Kur}, which is stated for an arbitrary discrete cocompact subgroup, torsion and inversions allowed; see also \cite[III, 5.4.3]{BC} and \cite[Proposition~4.2]{JL}.
\end{proof}

If $W$ contains some $w_m$ with $p \mid m$, then $\Gamma_W$ contains elements of odd valuation, which may invert edges of $I$; Proposition~\ref{prop:kur} applies to $\Gamma_W$ all the same.

\subsection{The graph}

\begin{proposition}\label{prop:classes}
There is an isomorphism of graphs $\Gamma_+\backslash I \cong \mathcal{G}$ which carries the action of $\Gamma_W/\Gamma_+ \cong W$ on $\Gamma_+\backslash I$ to the maps $w_m$ of Section~1. Under it, the stabiliser in $\Gamma_+$ of the edge $[I]$ has order half that of the unit group of the right order of $I$.
\end{proposition}

\begin{proof}
This is the argument of \cite[Propositions~4.4 and~4.6]{Rib}, for $D = pq$, and of \cite[\S1, proof of Proposition~1.4]{BD98}, which allows any discriminant of $B'$ and any level; neither computes the stabilisers in terms of ideals or treats the Atkin--Lehner involutions $w_m$ with $p \nmid m$, so we write it out. Since $[R : S] = p$, we have $S_\ell = R_\ell$ and hence $S[1/p] = R[1/p]$. Let $v_0$ be the vertex of $I$ fixed by $\GL_2(\Z_p)$, put $v_1 = \varpi^{-1}v_0$, and let $e_0$ be the edge from $v_0$ to $v_1$, oriented that way; its stabiliser in $\GL_2(\Qp)$ is $\Qp^\times S_p^\times$. Call a vertex even or odd according to the parity of its distance from $v_0$; every edge then has one even and one odd end, and we orient it from the even one.

We first identify the vertices. For $g \in \GL_2(\Qp)$ let $I_g$ be the left ideal of $R$ with $(I_g)_p = R_p\,g^{-1}$ and $(I_g)_\ell = R_\ell$ for $\ell \neq p$; it depends only on the vertex $g v_0$ up to a central factor, and $I_{\gamma g} = I_g\,\gamma^{-1}$ for $\gamma \in R[1/p]^\times$. Every left ideal class of $R$ contains some $I_g$: for any $J$, the localisation $J[1/p] = R[1/p]\,J$ is principal, equal to $R[1/p]\,\beta$ for some $\beta$ \cite[Theorem~28.5.5]{Voi}, and $J\beta^{-1}$ equals $R_\ell$ at every $\ell \neq p$; its $p$-part is $R_p h$ for some $h$, so $J\beta^{-1} = I_{h^{-1}}$. If $I_{g'} = I_g x$ with $x \in B'^\times$, then $x \in R_\ell^\times$ for all $\ell \neq p$, so $x \in R[1/p]^\times$, and comparing $p$-parts gives $g'v_0 = x^{-1}g v_0$. Hence $g v_0 \mapsto [I_g]$ is a bijection from $\Gamma\backslash\{\text{vertices}\}$ to the left ideal classes of $R$. Elements of $\Gamma_+$ preserve parity and those of $\Gamma - \Gamma_+$, which exist by \cite[Corollary~28.6.4]{Voi}, reverse it; so every $\Gamma$-orbit of vertices contains even and odd vertices, and an element of $\Gamma$ taking an even vertex to an even one lies in $\Gamma_+$. Hence the $\Gamma_+$-orbits of even vertices and of odd vertices each correspond to the left ideal classes of $R$, giving the two copies $([J], 1)$ and $([J], 2)$.

For the edges, in the same way, $g e_0 \mapsto [I^S_g]$, with $(I^S_g)_p = S_p\,g^{-1}$ and $(I^S_g)_\ell = S_\ell$ for $\ell \neq p$, is a bijection from the $\Gamma$-orbits of oriented edges to the left ideal classes of $S$. Since every edge has a unique orientation from even to odd, preserved by $\Gamma_+$, this is also a bijection from the $\Gamma_+$-orbits of edges. Take $g$ with $v_p(\det g)$ even, so that $g e_0$ is oriented from even to odd. Its ends are $g v_0$ and $g\varpi^{-1}v_0$, with ideals $I_g = R\,I^S_g$, since $R_pS_p = R_p$, and $I_{g\varpi^{-1}}$, whose $p$-part is $R_p\varpi g^{-1} = R_pS_p\varpi g^{-1}$, so that $I_{g\varpi^{-1}} = R\,J_p(S)\,I^S_g$. This is the incidence of $\mathcal{G}$.

Next, the stabilisers. An element of $\Gamma_+$ fixing the edge $g e_0$ fixes it pointwise, since it preserves parity. Its lifts to $R[1/p]^\times$, after removing a power of the central element $p$, lie in $g S_p^\times g^{-1}$ and in $S_\ell^\times$ for $\ell \neq p$, that is, in the unit group of the right order of $I^S_g$; these units have reduced norm $1$, since $B'$ is definite, and so lie in $\Gamma_+$. The centre $\pm p^{\Z}$ of $R[1/p]^\times$ meets this unit group in $\pm1$, so the stabiliser has order half that of the unit group of the right order of $I^S_g$.

Finally, the action of $W$. For $p \nmid m$, the element $\alpha_m$ of Lemma~\ref{lem:wm} has reduced norm of even valuation, so it preserves parity, and $I_{\alpha_m g} = J_m(R)\,I_g\,\alpha_m^{-1}$: the two sides agree at $p$, where $J_m(R)_p = R_p$; at $\ell \mid m$, where $J_m(R)_\ell = R_\ell\alpha_m$; and at the other $\ell$, where $\alpha_m$ is a unit. So $\alpha_m$ acts on vertices by $[J] \mapsto [J_m(R)\,J]$, and the same computation with $J_m(S)$ gives $[I] \mapsto [J_m(S)\,I]$ on edges. An element $\omega \in \Gamma - \Gamma_+$, which gives $w_p$ by Lemma~\ref{lem:wp}, sends $g v_0$ to $\omega g v_0$, with $I_{\omega g} = I_g\,\omega^{-1}$ in the same class and of the other parity, so it acts by $([J], i) \mapsto ([J], 3 - i)$. It sends the edge $g e_0$ to $\omega g e_0$, oriented from odd to even; with the even-to-odd orientation this is $\omega g\varpi^{-1}e_0$, whose ideal has $p$-part $S_p\varpi g^{-1}\omega^{-1}$, so that it is $J_p(S)\,I^S_g\,\omega^{-1}$ and $w_p$ acts on edges by $[I] \mapsto [J_p(S)\,I]$.
\end{proof}

\begin{proof}[Proof of Theorem~\ref{thm:main}]
By Lemma~\ref{lem:quot}, $Y$ has a projective model over $\Z_{p^2}$ with formal completion $\Gamma_W\backslash\Ohat$. By Proposition~\ref{prop:kur} this model is semistable, the components of its special fibre are rational curves with normalisations isomorphic to $\mathbb{P}^1$, and its dual graph is $(\Gamma_W\backslash I)^*$. Since $\Gamma_W/\Gamma_+ \cong W$, this is $(W\backslash(\Gamma_+\backslash I))^*$, which is $(W\backslash\mathcal{G})^*$ by Proposition~\ref{prop:classes}.
\end{proof}

\section{Application to hyperellipticity}\label{sec:application}

We recall the criterion proved in the paper (its theorem on dual graphs). Graphs may have loops and multiple edges; an automorphism permutes vertices and half-edges compatibly with incidence, so an involution may reverse a loop, and $\Gamma/\iota$ is the quotient of the topological realisation of $\Gamma$.

\begin{theoremA}
Let $X$ be a curve of genus $g \ge 3$ over a finite extension $K$ of $\Qp$. Suppose $X$ has a semistable model over $\mathcal{O}_K$ whose special fibre has only components of genus~$0$, and let $\Gamma$ be the core of its dual graph $\Gamma_0$. If $X$ is hyperelliptic over $\overline{K}$, then $\Gamma$ has an involution $\iota$ such that $\Gamma/\iota$ is a tree.
\end{theoremA}

Here the core of a graph is obtained by repeatedly deleting vertices of valence one, contracting bridges and suppressing vertices of valence two, a loop counting twice towards the valence. This is the form in which the proof of Theorem~A uses it: there the graph obtained from a subdivision of $\Gamma_0$ by contracting bridges is a subdivision of $\Gamma$, and an involution of it induces one of $\Gamma$ because it preserves the vertices of valence at least three.

\begin{corollary}\label{cor:test}
Let $D$, $N$, $W$, $p$ be as in Section~1, with $N$ not necessarily squarefree, and suppose $Y = X_0^D(N)/W$ has genus at least $3$. If the core of $(W\backslash\mathcal{G})^*$ has no involution with tree quotient, then $Y$ is not hyperelliptic.
\end{corollary}

\begin{proof}
By Theorem~\ref{thm:main}, $Y$ has a semistable model over $\Z_{p^2}$ whose special fibre components have normalisations isomorphic to $\mathbb{P}^1$ and whose dual graph is $(W\backslash\mathcal{G})^*$. Theorem~A, applied with $K = \Q_{p^2}$, gives the result. A component of the special fibre may meet itself, giving a loop in the dual graph; the proof of Theorem~A passes to a ramified extension and resolves the singularities, which turns such a loop into a cycle.
\end{proof}

\begin{thebibliography}{ABBR}

\bibitem[BC91]{BC}
J.-F.~Boutot, H.~Carayol,
\emph{Uniformisation $p$-adique des courbes de Shimura: les théorèmes de Čerednik et de Drinfeld},
Astérisque \textbf{196--197} (1991), 45--158.

\bibitem[BD98]{BD98}
M.~Bertolini, H.~Darmon,
\emph{Heegner points, $p$-adic $L$-functions, and the \v{C}erednik--Drinfeld uniformization},
Invent. Math. \textbf{131} (1998), 453--491.

\bibitem[BD96]{BD}
M.~Bertolini, H.~Darmon,
\emph{Heegner points on Mumford--Tate curves},
Invent. Math. \textbf{126} (1996), 413--456.

\bibitem[EGA]{EGA}
A.~Grothendieck,
\emph{\'El\'ements de g\'eom\'etrie alg\'ebrique III},
Publ. Math. IHES \textbf{11} (1961).

\bibitem[JL85]{JL}
B.~Jordan, R.~Livné,
\emph{Local diophantine properties of Shimura curves},
Math. Ann. \textbf{270} (1985), 235--248.

\bibitem[Kur79]{Kur}
A.~Kurihara,
\emph{On some examples of equations defining Shimura curves and the Mumford uniformization},
J. Fac. Sci. Univ. Tokyo Sect. IA Math. \textbf{25} (1979), 277--300.

\bibitem[Ogg85]{Ogg}
A.~Ogg,
\emph{Mauvaise réduction des courbes de Shimura},
Séminaire de Théorie des Nombres, Paris 1983--84,
Progr. Math. \textbf{59} (1985).

\bibitem[PS]{PS}
Padurariu, Saia,
arXiv:2509.25368.

\bibitem[Rib90]{Rib}
K.~Ribet,
\emph{On modular representations of $\mathrm{Gal}(\overline{\mathbb{Q}}/\mathbb{Q})$ arising from modular forms},
Invent. Math. \textbf{100} (1990), 431--476.

\bibitem[Voi21]{Voi}
J.~Voight,
\emph{Quaternion algebras},
Graduate Texts in Mathematics \textbf{288}, Springer, 2021. Numbering from the version on the author's homepage (v1.0.7u).

\end{thebibliography}

\end{document}
